\section*{Abstract.}

This software engineering project presents the design, development, and implementation of \textit{Vitalis PharmaSys}, a comprehensive web platform oriented towards the operational management and control of independent and community pharmaceutical establishments. The fundamental objective of the system is to systematize inventory administration, optimize medication dispensing, and ensure the sanitary traceability of products subject to medical prescription, mitigating economic losses caused by expiration and minimizing operational errors resulting from manual record-keeping. The technological solution is based on the Model-View-Controller (MVC) architectural pattern, developed using the Laravel framework on the PHP programming language in its 8.4 version, articulated with a MariaDB relational database management system under ACID transactions, and designed with Tailwind CSS in accordance with the \textit{Clinical Precision UI} design system guidelines. Its functional modules prominently include stock management through the FEFO (\textit{First Expired, First Out}) logistical strategy, automated emission of preventive alerts for expiration and stockouts, rigorous medical prescription control for antibiotics and psychotropics, an agile point-of-sale module, and an immutable operational audit log. The obtained results validate the system's effectiveness by drastically reducing counter service times, eliminating discrepancies in physical inventory, and ensuring strict compliance with applicable sanitary regulations, providing administrators with a reliable tool for strategic decision-making.

\vspace{0.5cm}
\noindent \textbf{Keywords:} Pharmacy Management System, Laravel Framework, Inventory Traceability, FEFO Logistics, Medical Prescriptions, MariaDB, Clinical UI.
